%%%%%%%%%%%%%%%%%%%%%%%%%%%%
% LaTeX poster template
% Created by Nathaniel Johnston
% August 2009
% http://www.nathanieljohnston.com/2009/08/latex-poster-template/
%%%%%%%%%%%%%%%%%%%%%%%%%%%%%%%%%%%%%%

\documentclass{beamer}
\usepackage[orientation=potrait,size=a0]{beamerposter}
\usepackage{graphicx}			% allows us to import images

%-----------------------------------------------------------
% The packages used in our latex file
%-----------------------------------------------------------

\usepackage[T1]{fontenc}
\usepackage{svg}
\usepackage{amsmath,amssymb,amsthm}
\usepackage{mathtools}
\mathtoolsset{centercolon}
\usepackage{tikz}
\makeatletter
\protected\def\tikz@nonactivecolon{\ifmmode\mathrel{\mathop\ordinarycolon}\else:\fi} 
\makeatother
\usetikzlibrary{decorations.pathreplacing}
%\usepackage{verbatim}
%\usepackage[normalem]{ulem} 
%\usepackage{thmtools}
%\usepackage{color}
%\usepackage[english]{babel}


%-----------------------------------------------------------
% Define the column width and poster size
% To set effective sepwid, onecolwid and twocolwid values, first choose how many columns you want and how much separation you want between columns
% The separation I chose is 0.024 and I want 4 columns
% Then set onecolwid to be (1-(4+1)*0.024)/4 = 0.22
% Set twocolwid to be 2*onecolwid + sepwid = 0.464
%-----------------------------------------------------------

\newlength{\sepwid}
\newlength{\onecolwid}
\newlength{\twocolwid}
\newlength{\threecolwid}
\setlength{\sepwid}{0.04\paperwidth}
\setlength{\onecolwid}{0.28\paperwidth}
\setlength{\twocolwid}{0.6\paperwidth}
\setlength{\threecolwid}{.92\paperwidth}
\setlength{\topmargin}{-0.5in}
\usetheme{confposter}
\usepackage{exscale}

%-----------------------------------------------------------
% The next part fixes a problem with figure numbering. Thanks Nishan!
% When including a figure in your poster, be sure that the commands are typed in the following order:
% \begin{figure}
% \includegraphics[...]{...}
% \caption{...}
% \end{figure}
% That is, put the \caption after the \includegraphics
%-----------------------------------------------------------

\usecaptiontemplate{
\small
\structure{\insertcaptionname~\insertcaptionnumber:}
\insertcaption}

%-----------------------------------------------------------
% Define colours (see beamerthemeconfposter.sty to change these colour definitions)
%-----------------------------------------------------------

\setbeamercolor{block title}{fg=ngreen,bg=white}
\setbeamercolor{block body}{fg=black,bg=white}
\setbeamercolor{block alerted title}{fg=white,bg=dblue!70}
\setbeamercolor{block alerted body}{fg=black,bg=dblue!10}

%%%%%%%%%%%%%%%%%%%%%%%%%%%%%%%%%%%%%%
% Miscellaneous Macros

\newcommand{\sbra}[1]{\langle{#1}|}
\newcommand{\sket}[1]{|{#1}\rangle}
\newcommand{\sbraket}[2]{\langle{#1}|{#2}\rangle}
\newcommand{\ceil}[1]{\lceil{#1}\rceil}
\newcommand{\CC}{\mathbb{C}}
\newcommand{\II}{\mathbb{I}}
\newcommand{\KK}{\mathbb{K}}
\newcommand{\QQ}{\mathbb{Q}}
\newcommand{\RR}{\mathbb{R}}
\newcommand{\ZZ}{\mathbb{Z}}
\newcommand{\posint}{\ZZ^+}
\renewcommand{\sc}{\mathrm{sc}}
\newcommand{\bc}{\mathrm{bc}}
\newcommand{\ph}{\mathrm{ph}}

\DeclareMathOperator{\spn}{span}


%-----------------------------------------------------------
% Name and authors of poster/paper/research
%-----------------------------------------------------------

\title{Universal computation by multiparticle quantum walk\\ with improved error bounds}
\author{Zak Webb} 

\institute{
  Institute for Quantum Computing, University of Waterloo\\
  Department of Physics \& Astronomy, University of Waterloo\\
  Department of Computer Science, University of Texas at Austin}
  
  
  \begin{document}

%----------------------------------------------------------
% Including the Logos
%----------------------------------------------------------

\addtobeamertemplate{headline}{} 
{
\begin{tikzpicture}[remember picture,overlay] 
  \node [shift={(12 cm,-5cm)}] at (current page.north west) {\includegraphics[height=10cm]{UW_Logo}};
\node [shift={(8.3cm,-10cm)}] at (current page.north west) {\includegraphics[height=5.66cm]{IQC_Logo}};
\node [shift={(-14.5 cm,-4.75cm)}] at (current page.north east) {\includegraphics[height=4cm]{UTCS_Logo}}; 

\end{tikzpicture} 
}


\tikzset{global scale/.style={
    scale=#1,
    every node/.style={scale=#1}
  }
}
%-----------------------------------------------------------
% Start the poster itself
%-----------------------------------------------------------

\begin{frame}[t]
\begin{columns}[t]
\begin{column}{\sepwid}\end{column}

\begin{column}{\onecolwid}
  \begin{alertblock}{Abstract}
    We prove that approximating the ground energy of the antiferromagnetic XY model on a simple graph at fixed magnetization is QMA-complete.  This strengthens a previous result considering a generalization of the XY model defined on graphs with self-loops.
  \end{alertblock}
  
  \begin{block}{XY model on a graph}
    The XY model is usually defined on a lattice, with $XX + YY$ terms between adjacent vertices.  The model on an arbitrary graph is the natural generalization of this.  Given a simple graph $G$ with vertex set $V(G)$ and edge set $E(G)$, the Hamiltonian is then
    \[
      O_G = \sum_{\{i,j\}\in E(G)} X_i X_j + Y_i Y_j.
    \]
  \end{block}
 
  
  \begin{block}{References}
  \nocite{*}
  \bibliographystyle{myhamsplain}
  \footnotesize{\bibliography{improved_bounds_poster}}
  \end{block}


\end{column}


\begin{column}{\sepwid}\end{column}

\end{columns}
\end{frame}



\end{document}